\section{Exact formal boundary and continuation}
The latest formally checked ancestor is the Block 49 source and its production, fresh extraction, and archive replay receipts. The relevant namespace for the physical operator is \texttt{KomorebiPhysics.Relativistic3DWeightedDomain}. The new proposed Lean file uses \texttt{Relativistic3DGeneratorBridge} and begins with exact finite cutoff commutation lemmas and physical Fourier transport. It intentionally does not declare the analytic target as a Lean theorem before supplying a genuine formal proof.

\subsection{Lean obligations still open}
\begin{enumerate}[leftmargin=1.5em,itemsep=.15em]
\item Establish the finite mode exponential derivative in the actual Hilbert spinor type and its coefficientwise estimate.
\item Pass that derivative to $\ell^2$ using summable graph weights, not a finite support assumption.
\item Prove the reverse strong generator domain inclusion using continuous coefficient evaluations.
\item Connect the resulting maximal generator with the self adjoint partial map and the spectral unitary group in the Lean library.
\item Define the everywhere bounded centered multiplier and prove its graph and resolvent convergence on the same coefficient space.
\item Add the third order Fourier weight and global $O(h^2)$ bounds with explicit source and target domains.
\end{enumerate}

\subsection{Physics and model obligations still open}
The background gauge field is not dynamic in this chapter. A lattice to continuum comparison requires an actual compatible family of finite grid Hilbert spaces, not only a centered symbol on $\ell^2(\Z^3)$. The gravity and backreaction monographs have separate constraint and physical interpretation requirements. The present analytic block neither resolves those questions nor changes the energy availability obstruction documented in the earlier backreaction study.

\section{Release ledger}
\begin{center}
\small
\begin{tabularx}{\linewidth}{>{\raggedright\arraybackslash}p{1.6in}X}
\toprule
\textbf{Component} & \textbf{Current status}\\\midrule
Block 49 source & Inherited, recorded checked v0.49.0 archive\\
Block 50 analytic theorems & Proved in the present mathematical manuscript; independent referee review pending\\
Block 50 finite computations & Executed in Python with recorded checks and one rejected negative control\\
Block 50 Lean draft & Source prepared, no compiler run in this environment\\
Strong lattice grid limit & Open, differs from centered Fourier multiplier limit\\
Dynamic gauge and gravity & Open and outside this chapter\\
\bottomrule
\end{tabularx}
\end{center}

\textbf{Restart point.} Reconstruct the pinned Lean environment from the Block 49 handoff, compile the finite cutoff bridge, and next formalize Theorem \ref{thm:generator} in its full domain form. Continue to the centered operator family after the generator is replayed. Preserve the earlier certified archive as an immutable predecessor while recording all repairs and fresh checks in a new Block 50 release ledger.

\subsection{Source concordance}
\begin{itemize}
\item \texttt{Relativistic3DFullFourierEvolution.lean}: unitary Fourier coefficient group, truncations, compact time convergence.
\item \texttt{Relativistic3DTorusFourierEvolution.lean}: full Fourier equivalence and physical torus propagation.
\item \texttt{Relativistic3DUnboundedDirac.lean}: maximal partial map, self adjointness, finite graph core.
\item \texttt{Relativistic3DWeightedDiracDomain.lean}: exact graph energy, mass independent first order domain, physical self adjoint Hamiltonian.
\item \texttt{Relativistic3DGeneratorBridge\_DRAFT.lean}: proposed next finite cutoff Lean interface, uncompiled.
\end{itemize}

\vspace{.25in}
\begin{center}\color{cyan}\rule{.28\linewidth}{.7pt}\end{center}
\vspace{.12in}
{\footnotesize\color{ash} \textit{Editorial note.} The analytic development is self-contained modulo the standard finite dimensional matrix exponential, dominated convergence for series, the spectral theorem for self adjoint operators, and elementary resolvent bounds. It is deliberately limited to the free torus model.}
